\makeabstract
    {
    Evaluating the Effect of RAS Inhibitors on T Cell Effector Functions using Mixed Lymphocyte Reactions (MLRs)
    }
    {
    Arum Han\textsuperscript{1}, Elisabeth Orgler-Gasche\textsuperscript{2,3}, Jiao Shen\textsuperscript{2,3}, Winnifer D. Conce Alberto\textsuperscript{2,3}, Michael Dougan\textsuperscript{3,4}, Stephanie Dougan\textsuperscript{2,3}
    }
    {
    \textsuperscript{1} Amherst College, Amherst, MA\\
\textsuperscript{2} Department of Cancer Immunology and Virology, Dana-Farber Cancer Institute, Boston, MA\\
\textsuperscript{3} Department of Immunology, Harvard Medical School, Boston, MA\\
\textsuperscript{4} Division of Gastroenterology, Massachusetts General Hospital, Boston, MA
    }
    {
    Mutations in RAS are among the most common oncogenes in cancer, contributing to approximately 20\% of all human cancers and more than 90\% of pancreatic ductal adenocarcinoma (PDAC) cases alone. Oncogenic RAS proteins have long been difficult to target, but recent advances in cancer therapies have led to the development of various RAS inhibitors. Although these agents are designed to suppress tumor cell signaling, their effects on immune cells, particularly T cells, are unknown. This study aims to investigate the effects of two RAS inhibitors, Daraxonrasib, a pan-RAS(ON) inhibitor, and MRTX1133, a selective KRAS G12D inhibitor, to compare how broad versus mutant-selective RAS inhibition influences T cell effector functions using in vitro mixed lymphocyte reactions (MLRs). 
    }
    {
    Mixed lymphocyte reactions (MLRs) were used as the experimental assay to evaluate the immunomodulatory effects of RAS inhibitors Daraxonrasib and MRTX1133. T cells from one donor are co-cultured with CD14+ antigen-presenting cells from another donor, and recognition of the foreign MHC by the TCR induces T cell activation. T cell responses following treatment with differing dosages of each RAS inhibitor were assessed by ELISA for IFN\ensuremath{\boldsymbol{\gamma}} secretion and flow cytometry for proliferation and activation marker expression.
    }
    {
    Overall, this study demonstrates that RAS inhibitors differentially modulate T cell effector functions in MLRs. Daraxonrasib, a pan-RAS(ON) inhibitor, produced broader suppression of activation, proliferation, and IFN\ensuremath{\gamma} production, whereas MRTX1133, a selective G12D KRAS inhibitor, inhibited T cell effector function at supraphysiologic concentrations. 
    }
    {
    These findings suggest that broad and mutant-selective RAS inhibitors have dose-dependent immunomodulatory effects that may influence anti-tumor immune responses. Future studies should further investigate how RAS inhibition affects specific T cell subsets, such as regulatory T cells. Evaluating changes in T cell populations and functions will provide a deeper understanding of how RAS inhibitors influence the tumor microenvironment.
    }

    \newpage
    